\scanpage{400}\mspage{139}
en pointillés $\alpha$ est la section de $\overline{X'}$ sur $X'$
déduite de la section $\tau$ de $\mathrm{Ch}_{\infty}(X)$ sur $X$,
qui réalise la décomposition de $f$ en
\[
X'\xrightarrow{\alpha}\overline{X'}\xrightarrow{\bar f}X
\]
où $\bar f$ devrait être une flèche ``\uline{fibrante}'' et
une $\mathrm{Hot}$-équivalence, $\alpha$ une flèche ``\uline{cofibrante}''
au sens de Thomason. (\deleted{\unsure{M}}Mais s'il en est ainsi,
est à voir !) Voici pour le moment les
propriétés essentielles que je voudrais vérifier.

$1^\circ$) $p$ est lisse et propre, et \inserted{est} une $\mathrm{HOT}$-fibration. \deleted{\unsure{De}}

$2^\circ$) Donc il en est de même de $\sigma=\mathrm{pr}_1p$ et $\beta=\mathrm{pr}_2p$.
De plus, ce sont des $\mathrm{HOT}$-fibrations \uline{triviales}, i.e.
à fibres \deleted{\unsure{contractiles}} \inserted{asphériques}.

$3^\circ$) Plus généralement \inserted{(pour ce qui concerne $\beta$)}, $\overline{X'}\xrightarrow{\bar f}X$ devrait \inserted{encore} être
propre et lisse et une $\mathrm{HOT}$-fibration, mais bien
sûr \deleted{\unsure{de}} \inserted{en général} pas triviale. (Elle devrait être triviale ssi
$f$ est une $\mathrm{Hot}$-équivalence.)
% slanted
\margin{\unsure{résulte} \unsure{formellement} de $1^\circ$ \deleted{\unsure{et $2^\circ$}} \unsure{et de} \unsure{cf p.|cf p.$^{\circ}$} 170}

Un mot sur la ``dualité''. Les énoncés sur
$\sigma$ (resp. $\beta$) se traduisent en énoncés sur
$\beta$ (resp. $\sigma$) par passage de $\underline t$ à $\underline t^{\circ}=(\deleted{\unsure{p}}t_n^{\circ})$,
où $\unsure{\tau|z}\unsure{\relbar|\longrightarrow}\tau^{\circ}$ est le foncteur de $\unsure{\mathrm{Ty}|\mathrm{Tz}}$ ``passage
à \deleted{\unsure{au dual}} l'intervalle opposé''. \unsure{On a|Oui|Or}, en effet
% archive sheet number: 399
